\documentclass[../main.tex]{subfiles}

\begin{document}

\section{Similarity transformations of operators: computing \texorpdfstring{$e^{A}Be^{-A}$}{eABe-A}}
\label{sec:similarity-transformations}

\subsection{Motivation}

Similarity transformations $B\mapsto e^{A}Be^{-A}$ are ubiquitous whenever a Hamiltonian is
brought to a simpler (often diagonal) form: Bogoliubov transformations of quadratic
Hamiltonians, displacement/squeezing of bosonic modes, polaron (Lang-Firsov)
transformations, Schrieffer-Wolff transformations, and passage to the interaction
picture or a rotating frame are all instances of this construction. If $A^\dagger=-A$
(generator anti-Hermitian), $U=e^{A}$ is unitary and the transformation preserves
Hermiticity and spectra. This is the case in all examples below.

\subsection{Method A: The Hadamard lemma}

\begin{lemma*}
For any two operators $A,B$,
\begin{align}
e^{A}Be^{-A} &= \sum_{n=0}^{\infty}\frac{1}{n!}\,\ad_A^n(B),\qquad \text{with} \\
\ad_A^n(B) &\equiv [A,B]_{n} \equiv [A, [A,B]_{n-1}] ,\quad \ad_A^0(B)\equiv[A,B] .
\end{align}
\end{lemma*}

This series is especially helpful whenever the nested commutator yields zero at for a finite $n$ or if the nested commutators eventually reproduce a small, recognizable pattern.

\paragraph{Strategy 1: Reduce to elementary operators.}
Utilize that for any $A,X,Y$,
\begin{equation}
e^{A}(XY)e^{-A}=\big(e^{A}Xe^{-A}\big)\big(e^{A}Ye^{-A}\big).
\end{equation}
It thus suffices to compute $e^{A}\,a\,e^{-A}$ and $e^{A}\,a^\dagger e^{-A}$ once and any
polynomial in these follows by substitution and normal-ordering, with no further
Hadamard series needed.

\paragraph{Strategy 2: The dagger shortcut.}
If $A^\dagger=-A$, $U=e^A$ is unitary with $U^\dagger=e^{-A}=U^{-1}$, so
$e^{A}Be^{-A}=UBU^\dagger$. Therefore, taking the dagger simply yields
\begin{equation}
\big(e^{A}Be^{-A}\big)^\dagger = e^{A}B^\dagger e^{-A}.
\end{equation}
Once $e^{A}ae^{-A}$ is known, $e^{A}a^\dagger e^{-A}$ follows for free by Hermitian
conjugation, roughly halving the bookkeeping.

\subsection{Method B: The generator flow method (not proof read since not used yet)}

Define $X(\theta)\equiv e^{\theta A}\,X\,e^{-\theta A}$, $X(0)=X$. Differentiating,
\begin{equation}
\frac{dX(\theta)}{d\theta}=[A,X(\theta)]=\ad_A\big(X(\theta)\big).
\end{equation}
This ordinary differential equation (ODE) is exactly equivalent to Method A -- the Hadamard series is its formal Taylor solution around $\theta=0$ -- but is can be more tractable in practice.

\paragraph{Key fact: closure $\Rightarrow$ finite-dimensional linear ODE.}
If $\ad_A$ maps a finite set of operators $\{X_1,\dots,X_n\}$ into linear combinations
of themselves, the operator ODE collapses to a matrix ODE:
\begin{equation}
\frac{d}{d\theta}\vec X(\theta)=M\,\vec X(\theta),\qquad
[A,X_j]=\sum_iM_{ij}X_i,\qquad \vec X(\theta)=e^{\theta M}\vec X(0),
\end{equation}
reducing the problem to exponentiating an ordinary $n\times n$ matrix.

% \paragraph{Why this works for Gaussian generators.}
% If $A$ is at most quadratic in bosonic ladder operators, $[A,a]$ and $[A,a^\dagger]$ are
% each at most linear, so $\{a,a^\dagger,\mathds{1}\}$ is automatically closed under
% $\ad_A$. Quadratic bosonic Hamiltonians generate the (metaplectic representation of the)
% symplectic group $Sp(2,\mathbb R)$ acting linearly on phase space.

% \subsubsection*{Worked examples}

% \paragraph{(a) Squeezing.} $A=\tfrac{\xi}{2}(a^2-a^{\dagger2})$: $[A,a]=\xi a^\dagger$,
% $[A,a^\dagger]=\xi a$, i.e.\ $M=\begin{pmatrix}0&1\\1&0\end{pmatrix}$. Eigenvalues
% $\pm1$ (hyperbolic), $M^2=\mathds{1}$:
% \begin{equation}
% a(\xi)=a\cosh\xi+a^\dagger\sinh\xi.
% \end{equation}

% \paragraph{(b) Displacement.} $A=\alpha(a^\dagger-a)$: $[A,a]=-\alpha\,\mathds{1}$ is a
% c-number, so enlarge the closed set to $\{a,\mathds{1}\}$. The matrix is nilpotent,
% giving polynomial growth:
% \begin{equation}
% a(\alpha)=a-\alpha\,\mathds{1}.
% \end{equation}

% \paragraph{(c) Phase rotation (interaction picture).} $A=i\theta\,a^\dagger a$:
% $[a^\dagger a,a]=-a$, closed on $\{a\}$ alone:
% \begin{equation}
% a(\theta)=e^{-i\theta}a.
% \end{equation}

% \paragraph{Classification.} Eigenvalues of $M$ purely imaginary $\to$ oscillatory
% (elliptic); real nonzero $\to$ hyperbolic (squeezing); nilpotent/zero $\to$ polynomial
% (displacement).

\subsection{When to prefer which method}
Both methods are mathematically identical. Method A is fast when the pattern is
recognizable by eye; Method B is more systematic for anything beyond the simplest
single-operator case, especially with mixed or multiple generators to track jointly.

\subsection{Scope: when this machinery fails}
The closure property requires $A$ to be at most quadratic in a set of canonically
(anti-) commuting operators. For a genuinely quartic operator such as
$H_{\mathrm{int}}=U\sum_i(n_i-\tfrac12)(n_{i+1}-\tfrac12)$, $\ad_A(H_{\mathrm{int}})$
does not close on a finite set for generic one-body $A$ (e.g.\ $A\propto J$ or
$A\propto T$) -- $[J,H_{\mathrm{int}}]$ is generically nonzero and does not reduce to a
small closed set of operators. The Hadamard series then cannot be resummed in closed
form; one is limited to a perturbative treatment in the strength of $A$.

\subsection{References}
\begin{itemize}
\item J.J.\ Sakurai \& J.\ Napolitano, \emph{Modern Quantum Mechanics}.
\item C.\ Gerry \& P.\ Knight, \emph{Introductory Quantum Optics}.
\item D.F.\ Walls \& G.J.\ Milburn, \emph{Quantum Optics}.
\item J.-P.\ Blaizot \& G.\ Ripka, \emph{Quantum Theory of Finite Systems}, MIT Press.
\item A.\ Perelomov, \emph{Generalized Coherent States and Their Applications}, Springer.
\item J.W.\ Negele \& H.\ Orland, \emph{Quantum Many-Particle Systems}.
\item B.C.\ Hall, \emph{Lie Groups, Lie Algebras, and Representations}.
\end{itemize}

\end{document}
